\chapter{Double Vision (Diplopia)}

\section*{Definition}
The perception of two images of a single object, classified as monocular (persists when one eye is covered, originates from ocular media) or binocular (resolves with covering either eye, originates from ocular misalignment).

\section*{What Happens in Your Body}
Binocular diplopia results from misalignment of the visual axes (strabismus) due to cranial nerve palsy (CN III, IV, VI), neuromuscular junction disorder (myasthenia gravis), restrictive orbitopathy (thyroid eye disease), or brainstem/cerebellar lesions affecting ocular motor control. Monocular diplopia arises from refractive errors, cataract, corneal irregularity, or retinal problems causing light scattering within one eye.

\section*{Causes}
\begin{itemize}
  \item Cranial nerve palsy (CN III, IV, VI from microvascular disease, aneurysm, trauma, tumor)
  \item Myasthenia gravis (fatigable diplopia, worse with sustained gaze)
  \item Thyroid eye disease (Graves orbitopathy, restrictive myopathy)
  \item Strabismus (childhood squint, decompensating in adulthood)
  \item Multiple sclerosis (demyelinating brainstem lesion)
  \item Stroke or transient ischemic attack (brainstem, cerebellum)
  \item Orbital trauma or fracture
  \item Cataract (monocular diplopia)
  \item Corneal irregularity (dry eye, keratoconus, scar)
  \item Botulism or Guillain-Barre syndrome (Miller Fisher variant)
\end{itemize}

\section*{When to See a Doctor}
See your doctor if:
\begin{itemize}
  \item Symptoms are severe or getting worse over time
  \item Sudden-onset diplopia (possible stroke, aneurysm, or brainstem lesion), diplopia with headache, ptosis, or pupil changes
  \item You have other concerning symptoms such as fever, weight loss, or bleeding
\end{itemize}

\section*{Related Symptoms}
\begin{itemize}
  \item Headache
  \item Ptosis (drooping eyelid)
  \item Blurred vision
  \item Dizziness
  \item Strabismus
\end{itemize}
\textbf{Important:} If this symptom persists, your doctor will check for serious underlying causes.

\section*{Self-Care / Home Management}
\begin{itemize}
  \item Rest and avoid activities that make it worse.
  \item Maintain adequate hydration and a balanced diet.
  \item Over-the-counter remedies may provide symptomatic relief (consult a pharmacist or doctor).
  \item Monitor symptoms and seek professional advice if they persist or worsen.
  \item Avoid self-medicating with prescription medications without medical guidance.
\end{itemize}

\section*{How Your Doctor Will Evaluate You}
\begin{itemize}
  \item A thorough history and physical examination are the first steps in evaluation.
  \item Laboratory tests (blood work, urinalysis) may be ordered based on clinical suspicion.
  \item Imaging studies (X-ray, ultrasound, CT, MRI) may be indicated when structural pathology is suspected.
  \item Specialist referral is arranged when the diagnosis remains uncertain or requires specific expertise.
\end{itemize}

\section*{Treatment Options}
\begin{itemize}
  \item Treatment targets the underlying cause identified through diagnostic evaluation.
  \item Supportive care and symptom management are provided while awaiting definitive therapy.
  \item Medications, lifestyle modifications, and physical therapies may be prescribed as appropriate.
  \item Follow-up ensures treatment effectiveness and allows timely adjustments to the care plan.
\end{itemize}

\clearpage